\chapter{Ciência, Saúde e Mundo Físico: da Previsão ao Benefício}
\label{ch:ciencia}

Prever uma estrutura molecular, sugerir uma hipótese e melhorar a vida de uma pessoa são resultados diferentes. A IA pode participar no caminho entre eles, mas não o percorre inteiro por decreto. A pergunta útil é onde acrescenta informação ou capacidade e que verificação continua por fazer.

\section{Uma aceleração científica com objeto definido}
AlphaFold apresentou resultados de previsão de estruturas de proteínas avaliados no CASP14, uma avaliação cega que permite comparar previsões com estruturas ainda não publicamente disponíveis~\citep{jumper2021}. É um avanço num problema concreto. Não equivale a explicar toda a biologia de uma proteína nem a demonstrar que um medicamento é eficaz e seguro.

Entre estrutura prevista e intervenção existem interações, dinâmica, contexto celular, hipóteses e experiências. Um sistema pode reduzir trabalho numa etapa e deslocar o gargalo para a seguinte. Esta é uma forma útil de procurar oportunidades: identifica a etapa lenta, o objeto que o modelo prevê e a experiência capaz de a contrariar.

Em investigação assistida, conserva hipóteses, dados e resultados separados. Uma sugestão de composto é uma proposta; uma simulação é um resultado de um modelo; uma medição em laboratório tem o seu procedimento e incerteza. Apagar estas distinções torna a investigação mais convincente no papel e menos reprodutível.

\section{Saúde: prever bem não é cuidar bem}
Um modelo pode ter bom desempenho num conjunto de imagens e falhar quando mudam equipamento, população ou prevalência. A utilidade clínica depende de decisões afetadas, falsos positivos, falsos negativos, integração e benefício observado. Uma avaliação de discriminação não substitui calibração, nem ambas substituem um estudo do processo de cuidado.

A orientação da OMS sobre grandes modelos multimodais discute benefícios possíveis e riscos como erros, enviesamento, privacidade e confiança excessiva na automação~\citep{who2024}. É orientação institucional; não certifica um produto. Este capítulo analisa critérios de avaliação, sem fornecer diagnóstico ou tratamento.

Uma tarefa administrativa, como resumir uma consulta para revisão profissional, também precisa de confirmar negações, doses, datas e atribuição de sintomas. Chamar «apenas documentação» ao texto não elimina consequências se esse documento for usado depois numa decisão.

\section{Meteorologia não é climatologia}
GraphCast foi apresentado como método de previsão meteorológica de médio prazo treinado em dados de reanálise, com avaliação sobre alvos definidos~\citep{graphcast2023}. Uma comparação desse tipo permite discutir previsões de variáveis e horizontes específicos. Não demonstra que o sistema resolve todas as escalas de previsão nem que substitui a física e a observação que contribuíram para os dados.

Para clima, eventos raros ou condições fora do histórico, o horizonte e o problema mudam. Pergunta que dados de referência existem, que distribuição se está a avaliar e como se tratam incerteza e extremos. «Previsão» é uma palavra demasiado ampla para transportar uma percentagem de um estudo para outro domínio.

\section{Indústria e agricultura: exemplos de desenho}
Num caso industrial conceptual, sensores e registos podem apoiar manutenção preditiva. O resultado esperado não deve ser apenas «deteta anomalias», mas uma decisão útil antes de uma falha, com custo de alarme falso e de intervenção. Se as máquinas mais antigas são inspecionadas mais vezes, os registos podem refletir a política de inspeção além da condição física.

Na agricultura, uma proposta de classificar sinais de doença em folhas exige testar iluminação, variedade, região, estação e doenças visualmente semelhantes. Um conjunto de fotografias de laboratório não representa automaticamente uma exploração agrícola. Trata estes exemplos como perguntas de projeto, não como resultados experimentais atribuídos a um produto.

Em ambos, os operadores conhecem exceções e limites que podem não constar dos dados. A recolha dessa informação é trabalho técnico. Instalar um modelo antes de perceber como o sensor falha é uma maneira rápida de automatizar uma interpretação errada.

\section{Robótica acrescenta consequências físicas}
Uma ação errada num chat pode ser revista antes de uso. Uma ação errada de um braço robótico pode colidir com um objeto ou uma pessoa. Perceção, planeamento, controlo e limites físicos precisam de avaliação conjunta. Demonstrações gravadas devem indicar condições, supervisão, tentativas e intervenções.

RT-2 estuda uma integração de visão, linguagem e ações representadas por tokens em tarefas robóticas~\citep{rt22023}. Isso mostra uma linha de investigação sobre transferência de representações. Não demonstra autonomia doméstica geral. A mesma instrução pode enfrentar geometria, força e objetos que nunca apareceram na avaliação.

Uma estratégia de engenharia separa proposta de plano e controlo físico com limites. A validação pode começar em simulação e avançar para ambientes controlados. A passagem ao mundo físico continua a exigir análise própria de segurança; não se resolve com uma frase «tem cuidado» no prompt.

\section{Benefício público e acesso}
Mesmo uma ferramenta eficaz pode não chegar a quem mais precisa se exigir equipamento, dados, conectividade ou formação indisponíveis. A avaliação deve incluir a população e as condições de uso pretendidas. Métricas agregadas podem esconder grupos para os quais o erro é maior.

Uma pergunta completa combina capacidade e distribuição: melhorou o resultado, para quem, com que recursos e com que possibilidade de contestação? É assim que um benchmark se liga a uma decisão de implementação, sem fingir que as duas coisas são iguais.

\begin{promptlabbox}{Do resultado do artigo à decisão de projeto}
\textbf{Objetivo e contexto:} ler uma fonte primária de ciência aplicada. Usa um dos artigos identificados neste capítulo e uma ficha de evidência.

\textbf{Procedimento:} identifica objeto previsto, dados, comparação, horizonte e métrica. Escreve uma utilização que o estudo suporta discutir e outra que ainda exige evidência. Não experimentes intervenções clínicas ou físicas a partir do exercício.

\textbf{Resultado esperado e validação:} cada conclusão tem localizador e alcance. A passagem da previsão ao benefício identifica as etapas que faltam.

\textbf{Falhas possíveis:} trocar acerto em benchmark por eficácia clínica, generalizar entre horizontes ou apresentar um caso conceptual de agricultura como estudo executado.
\end{promptlabbox}
